\documentclass[11pt,a4paper]{article}
% Generated by script/extract_concerns.py -- do not edit by hand.
% The preamble below is inherited verbatim from minor_report/main.tex so that
% every macro the extracted passages use is defined the same way here.
\makeatletter
\def\input@path{{/home/ambushee/wil_project/docs/report/minor_report/}}
\makeatother


\usepackage[T1]{fontenc}
\usepackage[utf8]{inputenc}
\usepackage{lmodern}
\usepackage{geometry}
\usepackage{microtype}
\usepackage{booktabs}
\usepackage{tabularx}
\usepackage{longtable}
\usepackage{array}
\usepackage{enumitem}
\usepackage{xcolor}
\usepackage{amsmath}
\usepackage{amssymb}
\usepackage{ifthen}
\usepackage{siunitx}
\usepackage{graphicx}
\usepackage{tikz}
\usetikzlibrary{arrows.meta,fit,positioning,shapes.geometric}
% sidewaystable: the per-run trajectory table is 18 columns wide and does not
% fit the portrait text block.
\usepackage{rotating}
\usepackage{hyperref}
\usepackage[nameinlink,noabbrev]{cleveref}

\geometry{margin=25mm}
\hypersetup{
  colorlinks=true,
  linkcolor=blue!45!black,
  urlcolor=blue!55!black,
  citecolor=blue!45!black,
  pdftitle={WiL Minor Progress Report},
  pdfauthor={WiL Project}
}
\setlist{nosep}
\setlength{\emergencystretch}{3em}

\definecolor{implemented}{HTML}{D9EAD3}
\definecolor{external}{HTML}{D9EAF7}
\definecolor{blocked}{HTML}{F4CCCC}
\definecolor{planned}{HTML}{FFF2CC}

% Per-block performance verdicts for the pipeline diagrams of section 4.9.
% These fills are GENERATED, never typed: figures/block_colours_<cond>.tex is
% written by script/block_performance.py from the run artifacts.
\definecolor{perfgood}{HTML}{A5D6A7}
\definecolor{perfpoor}{HTML}{FFCDD2}
% 'no claim' is yellow. That moved perfnoeffect off yellow and onto a blue tint:
% the two must stay visually distinct, because "measured, and it changes nothing"
% and "not measured at all" are opposite statements about the evidence.
\definecolor{perfnoeffect}{HTML}{D9EAF7}
\definecolor{perfnone}{HTML}{FFF2CC}
\newcommand{\perfmark}[1]{%
  \ifthenelse{\equal{#1}{good}}{\colorbox{perfgood}{\scriptsize good}}{%
  \ifthenelse{\equal{#1}{poor}}{\colorbox{perfpoor}{\scriptsize poor}}{%
  \ifthenelse{\equal{#1}{noeffect}}{\colorbox{perfnoeffect}{\scriptsize no effect}}{%
  \colorbox{perfnone}{\scriptsize no claim}}}}}

% Queried material. A query marks a result, number or condition that is measured
% but NOT yet trustworthy as evidence -- an unexplained mechanism, a confound, an
% unfair comparison, or a figure whose provenance does not support the use it is
% being put to. It is deliberately not used for "bad result": a large error that
% is correctly measured and understood is a finding, not a query.
%
% THE QUERIES NO LONGER RENDER IN THIS DOCUMENT. They are collected into the
% separate Open concerns document, which script/extract_concerns.py builds by
% PARSING THIS SOURCE. The markup therefore stays in the chapter files and is
% suppressed here rather than deleted -- deleting it would leave the concerns
% document with nothing to extract, and the chapter source is the single record
% of which passages are queried. To see them in place again, restore the two
% definitions marked RENDERING below.
\definecolor{flagred}{HTML}{B00000}

% \flag marks a queried VALUE: a table cell, a measured number, a status word,
% or a fragment inside a sentence. Those cannot be lifted out of the document --
% removing a table cell leaves a hole in the table, and removing a clause breaks
% the sentence around it -- so the value stays and only the colour is dropped.
\newcommand{\flag}[1]{#1}
% RENDERING: \newcommand{\flag}[1]{\textcolor{flagred}{#1}}

% \moved marks a queried passage that is a COMPLETE, self-contained sentence in
% running prose, so it can be lifted out without damaging what surrounds it. It
% renders as nothing here and appears only in the concerns document.
\newcommand{\moved}[1]{\unskip}

% Whole-paragraph query. Suppressed in this document; extracted in full by
% script/extract_concerns.py.
\newenvironment{queried}[1][Questionable]%
  {\setbox0\vbox\bgroup}%
  {\egroup}
% RENDERING:
% \newenvironment{queried}[1][Questionable]%
%   {\par\addvspace{0.6ex}\begingroup\color{flagred}\noindent
%    \textbf{[#1]}\ \ignorespaces}%
%   {\par\endgroup\addvspace{0.6ex}}

% Repository paths, identifiers, and commit hashes are typeset verbatim in a
% typewriter face, but must remain breakable: an unbreakable 40-character hash
% or a deep source path otherwise overruns the text block and prints on top of
% the neighbouring table column. Breaks are offered at path separators first
% (\allowbreak, penalty 0) and only reluctantly between other characters
% (penalty 4000), so a path divides after "/", "_" or "." whenever it can and
% splits a bare token such as a commit hash only when nothing else will fit.
% Hyphens are deliberately not separators: breaking "--symlink-install" after a
% hyphen would read as though a hyphen had been inserted.
\newcommand{\wilsoftbreak}{\penalty4000\relax}
\ExplSyntaxOn
\tl_new:N \l__wil_repo_tl
\NewDocumentCommand{\repo}{m}
  {
    \group_begin:
      \tl_set:Nx \l__wil_repo_tl { \tl_to_str:n {#1} }
      \regex_replace_all:nnN { ([a-zA-Z0-9]) } { \1 \c{wilsoftbreak} }
        \l__wil_repo_tl
      \regex_replace_all:nnN { ([/_.]) } { \1 \c{allowbreak} }
        \l__wil_repo_tl
      \texttt { \tl_use:N \l__wil_repo_tl }
    \group_end:
  }
\ExplSyntaxOff
\newcommand{\status}[1]{\textbf{#1}}

% Defines \pipelinefig and the block styles used by section 4.9. Loaded here
% rather than at the figure, because it declares \newcommand and \tikzset.
% Per-block performance diagram, drawn once and instantiated per condition.
%
%   \pipelinefig{static}    \pipelinefig{dynamic}
%
% Block fills come from figures/block_colours_<cond>.tex, which
% script/block_performance.py generates from the run artifacts. Nothing here
% chooses a colour; this file only decides where the boxes go. If a measurement
% changes, regenerate and the figure follows -- the alternative, typing ~40
% colours by hand, would put that many unsourced claims into a drawing that
% reads as authoritative.
%
% The structure follows the RY-SLAM framework figure in spirit, but the blocks
% are this project's: stereo-inertial rather than RGB-D, bounding boxes rather
% than segmentation masks, no motion gate, and no loop fusion inside VINS.
% Copying that paper's layout would assert capabilities this system lacks.

\newcommand{\bcol}[2]{\csname blockcolour@#1@#2\endcsname}

% Geometry note. The figure is scaled to \textwidth, so the ONLY thing that makes
% a connector visible is its length relative to the blocks either side. Wide blocks
% with narrow gaps -- 20mm against 2.2mm, as this was first drawn -- leave arrows
% about 2mm long after scaling, which reads as a smudge. Narrower blocks and much
% wider gaps cost nothing in total width and give the arrows four times the run.
\tikzset{
  pblock/.style={draw, rounded corners=1.5pt, minimum height=11mm,
    minimum width=18mm, align=center, font=\scriptsize, inner sep=2pt,
    line width=0.4pt},
  stagelab/.style={font=\scriptsize\bfseries, anchor=east, align=right},
  pipelab/.style={font=\small\bfseries, anchor=west},
  flowar/.style={-{Latex[length=2.4mm,width=1.8mm]}, gray!45!black,
    line width=0.7pt},
}

% #1 condition, #2 block id. The node name, position and label text follow at the
% call site, so this takes TWO arguments -- declaring three made it swallow the
% opening parenthesis of the node name and every block failed to parse.
\newcommand{\pb}[2]{\node[pblock, fill=\bcol{#1}{#2}] }

\newcommand{\pipelinefig}[1]{%
\begin{tikzpicture}[node distance=7mm and 8mm]
  % ---------------- ORB-SLAM3 ----------------
  \node[pipelab] (orbtitle) at (0,0) {ORB-SLAM3 (stereo-inertial)};
  \pb{#1}{orbextract} (o1) at (2.0,-1.4) {Extract ORB\\(+preproc, stereo)};
  \pb{#1}{orbimupre} (o2) [right=of o1] {IMU\\preintegration};
  \pb{#1}{orbposepred} (o3) [right=of o2] {Initial pose\\estimation};
  \pb{#1}{orbtracklocal} (o4) [right=of o3] {Track\\local map};
  \pb{#1}{orbkfdec} (o5) [right=of o4] {New keyframe\\decision};
  \pb{#1}{orbyolo} (o6) [right=of o5] {Dynamic feature\\removal (YOLO)};
  \node[stagelab] at ([xshift=-3mm]o1.west) {Tracking};
  \foreach \a/\b in {o1/o2,o2/o3,o3/o4,o4/o5} \draw[flowar] (\a) -- (\b);

  \pb{#1}{orbkfins} (m1) at (2.0,-3.6) {KeyFrame\\insertion};
  \pb{#1}{orbmpcull} (m2) [right=of m1] {MapPoints\\culling};
  \pb{#1}{orbnewpoints} (m3) [right=of m2] {New points\\creation};
  \pb{#1}{orblba} (m4) [right=of m3] {Local BA};
  \pb{#1}{orbkfcull} (m5) [right=of m4] {Local keyframes\\culling};
  \pb{#1}{orbimuref} (m6) [right=of m5] {IMU scale\\refinement};
  \node[stagelab] at ([xshift=-3mm]m1.west) {Local\\mapping};
  \foreach \a/\b in {m1/m2,m2/m3,m3/m4,m4/m5,m5/m6} \draw[flowar] (\a) -- (\b);
  \draw[flowar] (o5.south) -- ++(0,-4mm) -| (m1.north);

  \pb{#1}{orbplacerec} (l1) at (2.0,-5.8) {Place\\recognition};
  \pb{#1}{orbloopfusion} (l2) [right=of l1] {Loop fusion\\+ ess.\ graph};
  \pb{#1}{orbfullba} (l3) [right=of l2] {Full BA};
  \node[stagelab] at ([xshift=-3mm]l1.west) {Loop\\closing};
  \foreach \a/\b in {l1/l2,l2/l3} \draw[flowar] (\a) -- (\b);
  \draw[flowar] (m6.south) -- ++(0,-4mm) -| (l1.north);

  % ---------------- VINS-Fusion ----------------
  \node[pipelab] at (0,-7.5) {VINS-Fusion (stereo-inertial) + RTAB-Map};
  \pb{#1}{vinstrack} (v1) at (2.0,-8.9) {Feature\\tracking};
  \pb{#1}{vinsqueue} (v2) [right=of v1] {Feature\\queue};
  \pb{#1}{vinsyolo} (v3) [right=of v2] {Persistent-feature\\removal (YOLO)};
  \node[stagelab] at ([xshift=-3mm]v1.west) {Front end};
  \draw[flowar] (v1) -- (v2);

  \pb{#1}{vinsimuprop} (b1) at (2.0,-11.1) {IMU\\propagation};
  \pb{#1}{vinstri} (b2) [right=of b1] {Triangulation};
  \pb{#1}{vinsconstruct} (b3) [right=of b2] {Problem\\construction};
  \pb{#1}{vinssolve} (b4) [right=of b3] {Ceres solve};
  \pb{#1}{vinsmarg} (b5) [right=of b4] {Marginalization};
  \pb{#1}{vinsoutlier} (b6) [right=of b5] {Outlier\\rejection};
  \node[stagelab] at ([xshift=-3mm]b1.west) {Back end};
  \foreach \a/\b in {b1/b2,b2/b3,b3/b4,b4/b5,b5/b6} \draw[flowar] (\a) -- (\b);
  \draw[flowar] (v2.south) -- ++(0,-4mm) -| (b1.north);

  \pb{#1}{vinsslide} (c1) at (2.0,-13.3) {Slide\\window};
  \pb{#1}{vinssolverq} (c2) [right=of c1] {Solver\\convergence};
  \pb{#1}{vinsloop} (c3) [right=of c2] {Loop fusion\\(absent in fork)};
  \pb{#1}{vinsrtab} (c4) [right=of c3] {RTAB-Map\\correction};
  \node[stagelab] at ([xshift=-3mm]c1.west) {Output,\\back end};
  \draw[flowar] (c1) -- (c2);
  \draw[flowar] (b6.south) -- ++(0,-4mm) -| (c1.north);

  % ---------------- EKF, static only ----------------
  \ifthenelse{\equal{#1}{static}}{%
    \node[pipelab] at (0,-15.0) {EKF wheel\,+\,IMU (non-visual baseline)};
    \pb{#1}{ekfimupred} (e1) at (2.0,-16.4) {IMU\\prediction};
    \pb{#1}{ekfwheelupd} (e2) [right=of e1] {Wheel speed\\correction};
    \pb{#1}{ekfbiasobs} (e3) [right=of e2] {Gyro-bias\\observation};
    \pb{#1}{ekfyaw} (e4) [right=of e3] {Yaw-rate model\\(bicycle)};
    \pb{#1}{ekfcov} (e5) [right=of e4] {Covariance\\propagation};
    \pb{#1}{ekfanchor} (e6) [right=of e5] {Initial pose\\anchor};
    \node[stagelab] at ([xshift=-3mm]e1.west) {EKF};
    \foreach \a/\b in {e1/e2,e2/e3,e3/e4,e4/e5} \draw[flowar] (\a) -- (\b);
  }{%
    \node[font=\scriptsize\itshape, anchor=west, text width=0.82\linewidth]
      at (0,-15.4) {The EKF wheel\,+\,IMU baseline has no row here: no dynamic
      dataset carries wheel encoders or wheel odometry, so it cannot be run in
      this condition at all (\cref{tab:dataset-registry}).};
  }

  % ---------------- legend ----------------
  \begin{scope}[shift={(0,-18.4)}]
    \node[pblock, fill=perfgood, minimum height=8mm, minimum width=15mm] (g) at (2.2,0) {good};
    \node[font=\scriptsize, anchor=west] at ([xshift=1.5mm]g.east)
      {meets its criterion};
    \node[pblock, fill=perfpoor, minimum height=8mm, minimum width=15mm] (r) at (7.4,0) {poor};
    \node[font=\scriptsize, anchor=west] at ([xshift=1.5mm]r.east)
      {fails its criterion};
    \node[pblock, fill=perfnone, minimum height=8mm, minimum width=15mm] (n) at (12.4,0) {no claim};
    \node[font=\scriptsize, anchor=west] at ([xshift=1.5mm]n.east)
      {not instrumented, not exercised, or not implemented};
  \end{scope}
\end{tikzpicture}}






\graphicspath{{/home/ambushee/wil_project/docs/report/minor_report/}}

\geometry{margin=22mm}
\usepackage{titlesec}

% The inherited preamble SUPPRESSES queried material, because the report no
% longer renders it -- that is the whole point of this document existing. Undo
% the suppression here, or every passage extracted below would be typeset into
% nothing and this document would silently come out empty of the very text it
% was built to carry.
\renewcommand{\flag}[1]{\textcolor{flagred}{#1}}
\renewcommand{\moved}[1]{\textcolor{flagred}{#1}}
% Cross-reference targets live in the report, so a reference is printed, not
% resolved -- a dangling \cref would either fail or print a wrong number.
\newcommand{\reportref}[1]{\textcolor{gray}{[\texttt{\detokenize{#1}}]}}
\newcommand{\reportfigure}[1]{\textcolor{gray}{[figure: \texttt{\detokenize{#1}}]}}
\titlespacing*{\section}{0pt}{2.2ex plus .2ex}{0.8ex}
\titlespacing*{\subsection}{0pt}{1.6ex plus .2ex}{0.5ex}

\title{Open Concerns}
\author{Bolton Athit Davies}
\date{20 September 2026}

\begin{document}
\maketitle
\thispagestyle{empty}

\noindent
Every passage the minor progress report marks as \flag{queried}, collected into
one list. 77 items, from the minor progress report.

\medskip
\noindent
A flag marks a result, number or condition that is \emph{measured but not yet
trustworthy as evidence} --- an unexplained mechanism, a confound, an unfair
comparison, or a figure whose provenance does not support the use it is being put
to. It deliberately does not mark a bad result: a large error that is correctly
measured and understood is a finding, not a concern.

\medskip
\noindent
These passages remain in the report beside the claims they qualify, which is
where they belong. This document copies rather than moves them, so the whole set
can be read as a checklist. References such as \reportref{sec:example} point into
the report.

\medskip
\hrule

\section*{Chapter 3 --- Methodology}

\subsection*{Computing and software environment}
\noindent\textbf{1.}\quad \emph{table row:} OpenCV (Python) \textbar{} 4.11.0 \textbar{} \flag{differs from the 4.5.4 the estimators link}

\medskip

\subsection*{Simulated sensor rig}
\noindent\textbf{2.}\quad \emph{table row:} \texttt{acc\_n} \textbar{} 0.02 \textbar{} \flag{0.0139--0.0294} \textbar{} m\,s$^{-2}/\sqrt{\mathrm{Hz}}$

\medskip
\noindent\textbf{3.}\quad \emph{table row:} \texttt{gyr\_n} \textbar{} 0.002 \textbar{} \flag{0.000300--0.000364} \textbar{} rad\,s$^{-1}/\sqrt{\mathrm{Hz}}$

\medskip
\noindent\textbf{4.\ \flag{[Queried]}}\quad The configured gyroscope noise density is \num{2e-3} whereas the stationary
windows of the five recorded streams give \numrange{3.00e-4}{3.64e-4}, so the
configured value is about six times too large. The configured accelerometer
density of 0.02 sits inside the measured spread but that spread is itself wide,
\numrange{0.0139}{0.0294} across five bags of the same simulated sensor. Both
estimators are therefore running on an inertial noise model that does not
describe this IMU closely. The bias random walks cannot be checked at all: the
longest stationary interval in any \repo{dataset_allsensor} bag is
\SI{1.54}{\second} and an Allan-variance estimate needs minutes. No conclusion in
this report depends on these four numbers, but they should be re-derived from a
dedicated stationary recording before any tuning claim is made.

\medskip
\noindent\textbf{5.\ \flag{[Queried]}}\quad The accelerometer characterisation is not stable across recordings of the same
simulated sensor. \repo{allsensor_002}, whose stationary window is
\SI{1.54}{\second} and therefore the best conditioned of the three, returns an $x$
bias of \SI{0.0683}{\metre\per\second\squared}, against 0.0191 for
\repo{allsensor_000} and \SI{0.0016}{\metre\per\second\squared} for
\repo{allsensor_004}, whose window is the shortest at \SI{0.52}{\second}. The
scale ranges from 1.0199 to 1.0669 over the five. Because the window is short in
every case, part of this spread is estimator variance rather than a property of
the sensor, and the accelerometer terms should not be quoted as calibrated
constants. The gyroscope scale agrees to four decimal places across all five bags;
its bias does not, varying by a factor of twenty, which is consistent with the
same window-length problem.

\medskip

\subsection*{Simulation fidelity and its limits}
\noindent\textbf{6.}\quad \emph{table row:} Accelerometer, per axis (\si{\metre\per\second\squared}) \textbar{} \num{0.0566} \textbar{} \num{4.002e-3} \textbar{} \flag{\numrange{1.39e-2}{2.94e-2}} \textbar{} \num{2.00e-3}

\medskip
\noindent\textbf{7.}\quad \emph{table row:} \repo{slip1}, \repo{slip2} \textbar{} 0 on all grounds \textbar{} \flag{an ODE feature;
  \repo{dartsim} does not implement it, so inert at any value}

\medskip
\noindent\textbf{8.\ \flag{[Threat to validity]}}\quad This is a bias with a known direction. The EKF baseline of
\reportref{sec:results-proprio} is evaluated on a floor where the wheels effectively
cannot slip and with an encoder of unlimited resolution, which is dead
reckoning's best case. Both idealisations favour the EKF over the visual
estimators, and neither has been relaxed. The conclusion that a wheel+IMU filter
outperforms both visual pipelines should be read as holding \emph{under these
conditions}, not as a general result. \repo{script/make_slippery_floor.py}
generates ground models at a chosen friction coefficient so that traction can be
made an experimental factor; \repo{script/encoder_sim.py} imposes a realistic
quantisation floor. Neither has been used in any result reported here.

\medskip
\noindent\textbf{9.}\quad \emph{table row:} \repo{gyr_w} \textbar{} \num{0.0001} \textbar{} \num{1.735e-6} \textbar{} \flag{\num{58}$\times$} \textbar{} gyroscope bias random walk

\medskip

\subsection*{Real sensor rig and calibration limits}
\noindent\textbf{10.}\quad \emph{table row:} Second-camera intrinsics \textbar{} reuse of the first camera \textbar{} \flag{approximation}

\medskip
\noindent\textbf{11.}\quad \emph{table row:} IMU-to-camera translation \textbar{} zero \textbar{} \flag{placeholder}

\medskip
\noindent\textbf{12.}\quad \emph{table row:} Stationary $\lVert a\rVert$ \textbar{} \SI{9.2642}{\metre\per\second\squared} \textbar{} \flag{measured, deviates from \num{9.81}}

\medskip

\subsection*{Simulation environment}
\noindent\textbf{13.\ \flag{[Naming]}}\quad The factor recorded in this project as ``textured versus featureless floor'' is
narrower than the name suggests, and \reportref{fig:floor-treatments} is the evidence.
Both conditions rest on the \emph{same flat, untextured diffuse colour}; the stock
concrete-grain texture appears in neither. The entire difference is fourteen
painted bay outlines, \SI{0.145}{\metre} wide, covering a small fraction of a
\SI{1704}{\metre\squared} floor. A robot driving an aisle sees a uniform grey
surface in both conditions for most of its path, crossing a high-contrast strip
only at a bay boundary.

Results comparing the two floors must therefore be read as \emph{bay markings
versus nothing}, not as \emph{realistic floor versus nothing}, and the near-absence
of a visual-odometry improvement between them (\reportref{sec:results-proprio}) should
not be generalised to floors carrying genuine surface texture. Testing that would
require a world built on the stock \repo{GroundB_01} model, which no recorded
dataset uses.

\medskip
\noindent\textbf{14.\ \flag{[Unverified]}}\quad The dynamic world's baked grid, \reportref{fig:warehouse-maps}(c), contains
\SI{38}{\percent} fewer occupied cells than the distributed static arrangement
in~(a). \repo{bake_map.py} rasterises every \texttt{<include>} in the
world file at its authored pose, so kinematically driven props should appear at
their initial positions, but this has not been confirmed for the dynamic family.
Until it is, panel~(c) must not be used as occupancy truth for scoring a
reconstructed map of a dynamic scene: the deficit may be a genuine difference in
static stocking, or it may be props missing from the bake, and those two
possibilities imply opposite corrections. The estimator results in
\reportref{sec:results-proprio} are unaffected, as they use the arrangements in
panels~(a) and~(b) only.

\medskip

\subsection*{Dataset registry}
\noindent\textbf{15.}\quad \emph{table row:} \repo{dataset_static_nofloortexture_001} \textbar{} static \textbar{} 617 \textbar{} 1 \textbar{} 92.0 \textbar{} 113\,332 \textbar{} \checkmark \textbar{} \checkmark \textbar{} -- \textbar{} \flag{FD stub, WS ok}

\medskip
\noindent\textbf{16.}\quad \emph{table row:} \repo{dataset_dynamic_nofloortexture_01_000} \textbar{} dynamic \texttt{\_01} \textbar{} 713 \textbar{} 710 \textbar{} 165.4 \textbar{} 126\,790 \textbar{} \checkmark \textbar{} \checkmark \textbar{} -- \textbar{} \flag{FD malformed, WS ok}

\medskip
\noindent\textbf{17.}\quad \emph{table row:} \repo{dataset_real_001} \textbar{} real rig \textbar{} -- \textbar{} 2099 \textbar{} 85.5 \textbar{} 9\,629 \textbar{} -- \textbar{} -- \textbar{} -- \textbar{} \flag{0 IMU messages}

\medskip
\noindent\textbf{18.}\quad \emph{table row:} \repo{dataset_real_002} \textbar{} real rig \textbar{} -- \textbar{} 2869 \textbar{} 113.5 \textbar{} 12\,682 \textbar{} -- \textbar{} -- \textbar{} -- \textbar{} \flag{0 IMU messages}

\medskip
\noindent\textbf{19.\ \flag{[Data integrity]}}\quad Two drive copies are unreadable while presenting a complete, plausible
\repo{metadata.yaml}. \repo{dataset_static_nofloortexture_001} is a
\SI{524288}{\byte} stub whose metadata nevertheless declares 113\,332 messages,
and \repo{dataset_dynamic_nofloortexture_01_000} is \SI{710}{\mega\byte}---within
\SI{0.4}{\percent} of its workspace counterpart---with a valid SQLite header but a
store that reports \emph{database disk image is malformed} on any query. Both have
intact workspace copies, which is why those two datasets were read from the
workspace in the relogged campaign.

The general lesson is recorded as a control rather than an incident:
\repo{metadata.yaml} is written when recording begins and does not describe what
survived, so message counts, durations and topic lists copied from it are not
evidence that a bag is intact. Storage health must be established by opening the
SQLite store and counting rows, which is now part of the pre-experiment inspection
in \reportref{sec:collection-validity}.

\medskip
\noindent\textbf{20.\ \flag{[Scope limit]}}\quad Only the five \repo{dataset_allsensor} bags carry wheel encoders and wheel
odometry; the other thirteen recordings omit those topics entirely. The
proprioceptive baseline of \reportref{sec:results-proprio} therefore cannot be computed
for any dynamic condition, for the mapping bags, or for the real rig, and the
comparison between proprioception and visual odometry exists only in
\emph{static, featureless-floor} scenes---the conditions least favourable to
vision and most favourable to dead reckoning. Extending it requires re-recording
with the two topics bridged, not new analysis.

\medskip
\noindent\textbf{21.\ \flag{[Excluded]}}\quad \repo{dataset_real_001} and \repo{dataset_real_002} are excluded because their
inertial topic \repo{/hwt101ct_yaw_publisher} contains \emph{zero} messages, so
neither supports stereo-inertial operation at all. Their image streams are intact
(2\,562 and 3\,405 pairs), so they remain usable for detector or image-only work.
\repo{dataset_real_000} is the only real recording with inertial data, at 23\,257
messages.

\medskip

\subsection*{ORB-SLAM3 and Semantic Feature Filtering}
\noindent\textbf{22.}\quad \emph{table row:} \repo{loopClosing} \textbar{} 1 \textbar{} \flag{comment in the same file says ``set to 0
  deliberately''}

\medskip

\subsection*{Sensor model derived from data}
\noindent\textbf{23.}\quad \emph{table row:} accelerometer $x$ bias \textbar{} 0.0191 \textbar{} 0.0304 \textbar{} \flag{0.0683} \textbar{} 0.0309 \textbar{} 0.0016 \textbar{} \si{\metre\per\second\squared}

\medskip
\noindent\textbf{24.}\quad \emph{table row:} accelerometer $x$ scale \textbar{} 1.0199 \textbar{} 1.0249 \textbar{} \flag{1.0669} \textbar{} \flag{1.0525} \textbar{} \flag{1.0465} \textbar{} --

\medskip
\noindent\textbf{25.}\quad \emph{table row:} accelerometer $x$ scale \textbar{} 1.0199 \textbar{} 1.0249 \textbar{} \flag{1.0669} \textbar{} \flag{1.0525} \textbar{} \flag{1.0465} \textbar{} --

\medskip
\noindent\textbf{26.}\quad \emph{table row:} accelerometer $x$ scale \textbar{} 1.0199 \textbar{} 1.0249 \textbar{} \flag{1.0669} \textbar{} \flag{1.0525} \textbar{} \flag{1.0465} \textbar{} --

\medskip
\noindent\textbf{27.}\quad \emph{table row:} yaw-rate scale, bicycle \textbar{} 1.0110 \textbar{} 0.9527 \textbar{} 1.0034 \textbar{} \flag{0.9375} \textbar{} \flag{0.9250} \textbar{} --

\medskip
\noindent\textbf{28.}\quad \emph{table row:} yaw-rate scale, bicycle \textbar{} 1.0110 \textbar{} 0.9527 \textbar{} 1.0034 \textbar{} \flag{0.9375} \textbar{} \flag{0.9250} \textbar{} --

\medskip
\noindent\textbf{29.}\quad \emph{table row:} Window \textbar{} static duration \textbar{} 0.68 \textbar{} 0.58 \textbar{} 1.54 \textbar{} 1.04 \textbar{} \flag{0.52} \textbar{} \si{\second}

\medskip

\subsection*{EKF formulation}
\noindent\textbf{30.}\quad \moved{$\sigma_b$, the bias random walk, is assumed rather
than measured: the longest stationary window in any recording is
\SI{1.54}{\second} and an Allan-variance estimate needs minutes.}

\medskip

\subsection*{Reproducibility Controls and Implementation Deviations}
\noindent\textbf{31.}\quad \emph{table row:} Real-robot evaluation \textbar{} metric accuracy against ground truth \textbar{} calibration-limited, no independent trajectory truth \textbar{} \flag{blocked}

\medskip
\noindent\textbf{32.}\quad \emph{table row:} Metric infrastructure \textbar{} resource, throughput, uncertainty, manifests \textbar{} partially implemented; RPE now complete \textbar{} \flag{incomplete}

\medskip

\section*{Chapter 4 --- Results}

\subsection*{Question, hypotheses, and conditions}
\noindent\textbf{1.}\quad \emph{table row:} IMU noise model \textbar{} \flag{not controlled in this experiment.} The EKF's is measured
  per bag by \repo{proprio_estimator.py calibrate}; both visual estimators read
  four constants from \repo{config/wil_sim/stereo_imu.yaml}. The factor is
  confounded with the estimator here and is separated in
  \reportref{sec:proprio-tuning}

\medskip

\subsection*{Validity}
\noindent\textbf{2.}\quad \flag{The initialization and reset figures below were extracted for
\repo{allsensor_000} and \repo{allsensor_001} only and have not been recovered for
the other three bags:} VINS-Fusion initialized once in \SI{0.33}{\second} with no
resets, and both ORB-SLAM3 runs shut down gracefully with no tracking-loss events
and inertial initialization at \SI{2.48}{\second} and \SI{2.51}{\second}.

\medskip
\noindent\textbf{3.\ \flag{[Queried]}}\quad ORB-SLAM3 receives fewer stereo pairs than are available on every dataset, while
VINS-Fusion receives all of them (2639, 2478 of 2479, 2628, 3102 and 3046). The
ORB shortfall is \SIrange{3.8}{11.9}{\percent} and varies by run as well as by
dataset: \SI{90.2}{\percent} and \SI{94.2}{\percent} of the available pairs on
\repo{allsensor_000} for the two configurations, \SI{96.2}{\percent} and
\SI{95.1}{\percent} on \repo{allsensor_003}. Its own counters report zero decode
failures, and the stereo timestamps are aligned to \SI{0}{\second} median offset
with \SI{100}{\percent} of pairs inside the \SI{20}{\milli\second} synchronisation
slop, so synchronisation is not what loses them.

The counter itself cannot say where they go: \repo{stereo_pairs_received} is
incremented inside the \emph{synchronised} callback, and the node keeps no count of
images arriving. What the source does show is an asymmetry in input buffering.
ORB-SLAM3 subscribes with \repo{rclcpp::SensorDataQoS()}, which is best-effort with
a depth of five, so a frame arriving while the tracking callback is busy is
discarded by the middleware and never recorded anywhere; VINS-Fusion subscribes
with \repo{KeepLast(100).reliable()} and buffers the same burst. \moved{The shortfall
is therefore most consistent with ORB-SLAM3 dropping frames at its own subscriber
under load, not with loss in transport.} The loop-closure-free runs support that
reading on three datasets, receiving more pairs than the enabled runs on
\repo{allsensor_000}, \repo{_001} and \repo{_002} while doing less work per frame,
though \repo{allsensor_003} runs the other way and the mechanism is not proven.
Trajectory accuracy over a shared interval is not obviously affected---ORB-SLAM3
is the more accurate of the two on some datasets despite the shortfall---but no
throughput, latency or per-frame compute comparison between the two estimators is
admissible until the loss is explained: the two estimators differ not only in
load but in input-queue policy, so such a comparison would measure the QoS choice
as much as the algorithm.

\medskip
\noindent\textbf{4.\ \flag{[Ruled out]}}\quad A plausible reading of these results is that the visual estimators lose to the EKF
because their input rate collapsed. The data does not support it. Source imagery
is regular at \SI{30.3}{\hertz} in simulated time on all five bags, with a median
inter-frame interval of \SI{0.0330}{\second} and no gap exceeding
\SI{0.1}{\second}; VINS-Fusion received and processed every frame on every
dataset, with zero skipped before the back end, and still trails the
cross-validated EKF on four of five. Frame loss is real for ORB-SLAM3 and remains
unexplained, but it cannot account for the margin, because the estimator with no
frame loss at all shows the same one.

\medskip
\noindent\textbf{5.\ \flag{[Queried]}}\quad ORB-SLAM3 was run in two configurations on every dataset, recorded as
\repo{logging_20260916_01} with loop closing enabled and
\repo{logging_20260916_02} with it disabled, and both are reported below as
separate systems. Within the enabled configuration the loop closer behaves very
differently per dataset (\reportref{tab:orb-loop}), so those five trajectories are not
repeats of one condition and must not be pooled.

A counter reading \texttt{loop\_closures,0} does \emph{not} mean the loop closer
was inactive. On \repo{allsensor_000} it performed 447 place-recognition checks,
found 90 loop candidates and rejected all 90 on inertial geometry; on
\repo{allsensor_004} it found 114 and accepted none. The disabled runs are
distinguishable only in \repo{loop_closing.csv}, where every event carries
\texttt{current\_keyframe\_id}~$=-1$ and a duration below
\SI{0.003}{\milli\second}, the signature of a thread spinning with nothing queued
to it.

\medskip
\noindent\textbf{6.\ \flag{[Provenance gap]}}\quad \repo{run_metadata.csv} records 22 parameters and loop-closure state is not among
them, so the whole distinction rests on the directory name---and that name does
not mean the same thing everywhere in this chapter. In the relogged campaign of
\reportref{sec:results-validity}, \repo{logging_20260916_01} through \repo{_03} are
three \emph{repeats} of one configuration, all with an active loop closer. In the
\repo{allsensor} campaign reported here, \repo{_01} and \repo{_02} are two
\emph{configurations}, loop closing enabled and disabled, with one run each. The
same suffix therefore encodes a repeat index in one section and an experimental
factor in the next. The metadata for the two bag-000 runs is identical apart from the experiment
identifier, output path and start time, and both name the same installed
configuration file, so nothing in the recorded provenance distinguishes an
enabled run from a disabled one. The distinction used in this chapter was
recovered from \repo{loop_closing.csv} after the fact. The logger should record
the loop-closing and global-BA switches explicitly before any further ORB-SLAM3
campaign.

\medskip

\subsection*{Accuracy results}
\noindent\textbf{7.}\quad \emph{table row:} plain \textbar{} Wheel odometry \textbar{} 1.485 \textbar{} 2.888 \textbar{} 4.332 \textbar{} 3.76 \textbar{} \flag{$-30.11$}

\medskip
\noindent\textbf{8.}\quad \emph{table row:} \repo{allsensor_003} \textbar{} \texttt{bayline\_} \textbar{} IMU dead reckoning \textbar{} 45.562 \textbar{} 62.903 \textbar{} 74.745 \textbar{} \flag{53.16} \textbar{} $-5.84$

\medskip
\noindent\textbf{9.}\quad \emph{table row:} \texttt{objonwallonly} \textbar{} Wheel odometry \textbar{} 3.192 \textbar{} 3.979 \textbar{} 9.231 \textbar{} 6.57 \textbar{} \flag{$-37.17$}

\medskip
\noindent\textbf{10.}\quad \emph{table row:} ORB-SLAM3 \textbar{} \flag{6.481} \textbar{} -- \textbar{} -- \textbar{} -- \textbar{} --

\medskip
\noindent\textbf{11.}\quad \moved{That earlier figure
of ``3.0, 1.2 and 2.2 times'' was computed against ORB-SLAM3 with loop closing
enabled, which is not the best visual result available.} Measured against the best
visual estimator of each dataset, and using the cross-validated EKF
(\reportref{tab:proprio-xval}), the margins are 1.38, 1.41, 2.53, 1.31 and
\num{1.13}.

\medskip
\noindent\textbf{12.}\quad \flag{Nor did adding the bay markings help it} ---the two
\texttt{bayline} datasets give VINS-Fusion 0.784 and \SI{0.425}{\metre} against
0.616, 0.383 and \SI{0.558}{\metre} on the plain floors, which is no improvement.

\medskip

\subsection*{Falsification of H1, and the drift mechanism}
\noindent\textbf{13.}\quad H1 is \flag{falsified} across all five datasets.

\medskip
\noindent\textbf{14.}\quad \emph{table row:} \repo{allsensor_002} \textbar{} $+0.34$ \textbar{} \flag{$-30.11$} \textbar{} 865 \textbar{} $-3.481$ \textbar{} $-31.03$

\medskip
\noindent\textbf{15.}\quad \emph{table row:} \repo{allsensor_003} \textbar{} $-6.25$ \textbar{} \flag{$-37.17$} \textbar{} 1123 \textbar{} $-3.309$ \textbar{} $-38.51$

\medskip
\noindent\textbf{16.\ \flag{[Corrected]}}\quad An earlier reading of this experiment proposed that the offset came from the
bicycle model averaging the two kingpin angles, since $\tan$ of a mean is not the
mean of $\tan$. That explanation is \emph{wrong} and is withdrawn. Recomputing the
yaw rate as the mean of each wheel's own bicycle-equivalent rate changes the
residual from \SI{-0.00277}{} to \SI{-0.00288}{}, from \SI{-0.00285}{} to
\SI{-0.00296}{}, and from \SI{-0.00624}{} to \SI{-0.00630}{\radian\per\second}---
under \SI{4}{\percent} in every case, two orders of magnitude short of what the
observed drift requires. The averaging is not the source.

\medskip
\noindent\textbf{17.}\quad \emph{table row:} \repo{allsensor_003} \textbar{} bayline \textbar{} 0.9375 \textbar{} 0.1192 \textbar{} \flag{3.54} \textbar{} 0.423

\medskip
\noindent\textbf{18.}\quad \emph{table row:} \repo{allsensor_004} \textbar{} bayline \textbar{} 0.9250 \textbar{} 0.0613 \textbar{} \flag{3.41} \textbar{} 0.494

\medskip
\noindent\textbf{19.}\quad \flag{That inference
does not survive the textured-floor datasets} and is withdrawn: the coupling is
the dominant term in the EKF's heading error whenever the wheel rate is poor, and
it sets a floor on the filter's accuracy that no amount of gyroscope quality can
lift.

\medskip

\subsection*{Covariance consistency}
\noindent\textbf{20.}\quad \emph{table row:} IMU dead reckoning \textbar{} 311.0 \textbar{} 0.005 \textbar{} 0.022 \textbar{} \flag{0.043} \textbar{} overconfident

\medskip
\noindent\textbf{21.}\quad \emph{table row:} IMU dead reckoning \textbar{} 339.9 \textbar{} 0.001 \textbar{} 0.009 \textbar{} \flag{0.021} \textbar{} overconfident

\medskip
\noindent\textbf{22.}\quad \emph{table row:} IMU dead reckoning \textbar{} 244.7 \textbar{} 0.017 \textbar{} 0.045 \textbar{} \flag{0.147} \textbar{} overconfident

\medskip
\noindent\textbf{23.}\quad \emph{table row:} IMU dead reckoning \textbar{} \flag{1504.3} \textbar{} 0.010 \textbar{} 0.018 \textbar{} \flag{0.035} \textbar{} overconfident

\medskip
\noindent\textbf{24.}\quad \emph{table row:} IMU dead reckoning \textbar{} \flag{1504.3} \textbar{} 0.010 \textbar{} 0.018 \textbar{} \flag{0.035} \textbar{} overconfident

\medskip
\noindent\textbf{25.}\quad \emph{table row:} IMU dead reckoning \textbar{} 289.5 \textbar{} 0.004 \textbar{} 0.017 \textbar{} \flag{0.044} \textbar{} overconfident

\medskip
\noindent\textbf{26.\ \flag{[Method caveat]}}\quad The $\chi^2$ interval behind the verdict column assumes independent samples. At
\SI{200}{\hertz} consecutive dead-reckoning errors are almost perfectly
correlated, so the test is run on a \SI{1}{\hertz} subsample (81 to 102 samples per dataset).
That reduces the dependence but does not remove it, and the verdicts should be
read as indicative. The $\sigma$-coverage columns make no independence assumption
and are the sturdier statement.

\medskip

\subsection*{Cross-validation of the EKF noise model}
\noindent\textbf{27.\ \flag{[What this control does not cover]}}\quad Exchanging noise models between bags tests whether the filter saw \emph{this}
recording's answer. It cannot test the asymmetry that matters here, because both
models are still measured from data: all five bags come from one simulated IMU, so
the exchanged $\sigma$ values differ by less than a factor of two on every term
except accelerometer $x$. The visual estimators receive nothing of the kind---they
read four constants from a file. The control in \reportref{sec:proprio-tuning} replaces
this one.

\medskip

\subsection*{Transferring the visual systems' noise model to the EKF}
\noindent\textbf{28.\ \flag{[What the control does and does not establish]}}\quad \flag{H4 was written down after the exploratory runs had been scored, not
before.} What is reported here is a post-hoc formalisation of an exploratory
finding; the decision rule and the direction of the test were fixed before this
subsection was written, but not before the data was seen. The unanimity---every
bag, both directions, a mechanism that isolates to one term---is what makes it
worth reporting at this stage, not a pre-registered test, and a confirmatory
repeat on bags not used to form the hypothesis is the honest follow-up.

This is also an equal-handicap comparison, not an equal-effort one: both families
now run on a declared IMU noise that nobody fitted to the data. It establishes
that the EKF's advantage in \reportref{tab:proprio-accuracy} is smaller than the
tuning difference between the two families, and reverses under it. It does
\emph{not} establish that the visual systems would win a symmetric comparison.
The mirror experiment---giving VINS-Fusion and ORB-SLAM3 the measured
values---requires re-recording and has not been run to a reportable standard; the
one VINS run whose metadata confirms it diverged (\reportref{sec:sim-fidelity}). The
transfer is not perfectly symmetric in role either: the filter's $\sigma_{g}$ is
a scalar in a planar five-state model, whereas \repo{gyr_n} feeds three-axis
pre-integration. Finally, $n=1$ per cell is a property of a deterministic filter
and not a design choice, but it means the between-bag range is the only
uncertainty this experiment can report.

\medskip

\subsection*{Interpretation and limitations}
\noindent\textbf{29.\ \flag{[Generalisation limit]}}\quad The comparison is confined to \emph{static} scenes. Four worlds and both floor
conditions are now covered and the ranking is stable across them, so the earlier
form of this limitation---that the result held only on featureless floors, vision's
worst case---is no longer the binding one. \flag{Adding the bay markings did not improve
either visual estimator}, which removes that explanation rather than confirming
it---though the markings are a much weaker manipulation than ``textured floor''
suggests (\reportref{fig:floor-treatments}), so a genuinely textured surface remains
untested.

What remains is that every dataset carrying wheel encoders is static
(\reportref{tab:dataset-registry}). The result must not be restated as
``proprioception outperforms visual SLAM'' without that qualifier, and it cannot
be extended to dynamic scenes---where wheel odometry is unaffected by moving
objects, and where the visual systems' known failures occur---until the
\repo{allsensor} topics are recorded in a dynamic world. That recording remains
the single most informative next run in this line of work.

\medskip
\noindent\textbf{30.\ \flag{[Queried]}}\quad Both visual estimators replayed their bags at $0.5\times$ real time, so each frame
received roughly twice the wall-clock budget it would have online. The accuracy
figures in \reportref{tab:proprio-accuracy} are therefore not evidence of real-time
capability, and no timing conclusion may be drawn from this experiment. The three
baselines, being offline, have no replay rate at all, which is a second asymmetry
in the same comparison.

\medskip
\noindent\textbf{31.\ \flag{[Statistical limit]}}\quad $n=5$ across four worlds, and no condition is repeated more than twice. The tables
therefore report individual values deliberately; no mean, standard deviation or
confidence interval is computed, and between-dataset differences are not evidence
of a systematic effect. A single run per world cannot establish a scene-content
effect; it shows only that the estimator ranking survives changes of scene.

\medskip
\noindent\textbf{32.\ \flag{[Unmeasured variance]}}\quad ORB-SLAM3 is multi-threaded and not deterministic, and its run-to-run spread on
identical input has not been measured here: each of its two configurations was run
once per dataset. Several differences this chapter reports sit within a range that
spread could plausibly cover---most sharply \SI{6.481}{\metre} with loop closing
enabled on \repo{allsensor_003} against \SI{0.515}{\metre} with it disabled, and
the reversal on \repo{allsensor_002} where disabling it is three times
\emph{worse}. Those differences are real measurements of the runs that were made,
but they cannot yet be attributed to the loop-closing switch. Three or more
repeats per configuration, reported as median and spread in the manner of the
relogged campaign (\reportref{sec:results-validity}), are required before any
loop-closure conclusion is drawn.

\medskip

\subsection*{Trajectory Accuracy, Plausibility, and Robustness}
\noindent\textbf{33.\ \flag{[Excluded: dynamic-scene trajectory results]}}\quad Every dynamic-scene trajectory result in this section, and the anomalies A3 and
A4 that derive from them, are \textbf{withheld as findings} and retained only as
excluded evidence. The reason is the environment, not the measurement: the
dynamic worlds move their props along scripted kinematic paths. The props carry
collision geometry but produce no dynamic response against the robot, nothing
constrains their speeds or routes to warehouse traffic, and they can occupy the
same space as the vehicle without consequence. Numbers measured there describe
the simulation, not the problem the thesis poses.

The runs happened, the artifacts are intact, and the plausibility failures are
real properties of those recordings, so they are reported below rather than
deleted --- an excluded measurement that explains a methodological correction is
worth more in the record than a gap. What must not be done is to read them as
evidence about estimator behaviour in a dynamic warehouse.

Two consequences follow. The condition this work is named for currently has no
valid testbed: it is untested, not tested and failed. And the environment becomes
the leading candidate explanation for the breakdown itself, displacing the
open question left in \reportref{sec:anomaly-catalogue} --- testable by rebuilding the
worlds with physically simulated props and re-running, which requires no change
to either estimator.

Unaffected: the static-scene results here, the estimator comparison of
\reportref{sec:results-proprio}, the detector evaluation of
\reportref{sec:yolo-performance} --- prop motion is kinematic, which changes contact
physics rather than appearance, so detectability and class assignment are
untouched --- and the map evaluation of \reportref{sec:map-quality}.

\medskip

\subsection*{Interpretation and limitations}
\noindent\textbf{34.\ \flag{[Queried: reproducibility of the drivability half]}}\quad The geometry result is fully reproducible: both \repo{map/map_tiny.db} and the
world file are present. The drivability result is not. Its reference run,
\repo{dataset_static_tiny_repeat}, is absent from \repo{dataset/} and from the
removable drive, and the only surviving recording of this warehouse is the one
that built the map, which would answer a weaker question
(\reportref{sec:map-drivability}). The database additionally stores an empty parameter
set, so the configuration that produced it is not recoverable from the file, and
the RTAB-Map engineering record's claim that the node stamps match no dataset is
contradicted by the pose agreement reported above. The drivability figures should
therefore be read as a dated measurement rather than a standing result, and should
not be carried into the major report until an independent route has been recorded
and retained.

\medskip

\subsection*{Where each pipeline is weak, block by block}
\noindent\textbf{35.\ \flag{[Withdrawn criterion]}}\quad An earlier revision of these tables graded VINS-Fusion's solver quality against
``median final cost $>5000$''. That threshold was chosen after observing 1176 in
static scenes and 14110 in dynamic ones, so it separated two numbers already in
hand and the \flag{``poor'' verdict it produced was guaranteed by the choice
rather than tested by it}. Raw Ceres cost also has no absolute scale, depending on
residual count and weighting, so no fixed threshold on it is meaningful in
principle. It is replaced by the solver's own \repo{solver_solution_usable} flag,
which is binary and emitted by the solver rather than invented here:
\SI{99.64}{\percent} of solves in static scenes and \SI{99.55}{\percent} in
dynamic ones return a usable solution, and the block is now graded \perfmark{good}
in both. The non-convergence rate---\SI{22.9}{\percent} of solves in static scenes
and \SI{45.3}{\percent} in dynamic ones reaching the iteration or
\SI{80}{\milli\second} cap---is reported in the note column and deliberately not
graded, because any threshold on it would repeat the same error.

\medskip
\noindent\textbf{36.}\quad \textbf{Both visual front ends are the dominant cost, and neither keeps up
cleanly.} ORB-SLAM3 spends \SI{33.1}{\milli\second} of a \SI{53.8}{\milli\second}
tracking budget on preprocessing, extraction and stereo matching in static scenes,
rising to \SI{45.1}{\milli\second} of \SI{61.2}{\milli\second} in dynamic ones.
VINS-Fusion's back end waits \SI{403.7}{\milli\second} per frame on its feature
queue in static scenes and \SI{502.1}{\milli\second} in dynamic ones, against a
\SI{66.0}{\milli\second} frame budget, with a backlog of four to seven frames. \flag{These are the same deficit wearing different clothes} : neither estimator
processes frames as fast as they arrive, and the two differ in what they do about
it.

\medskip
\noindent\textbf{37.}\quad \emph{table row:} E \textbar{} \flag{replay rate $0.5\times$ for both visual estimators} \textbar{} accuracy figures
  are not evidence of real-time capability

\medskip
\noindent\textbf{38.}\quad \emph{table row:} E \textbar{} \flag{baselines are offline, visual estimators ran live} \textbar{} timing-induced
  error is present on one side of the comparison only

\medskip
\noindent\textbf{39.}\quad \emph{table row:} E \textbar{} \flag{only \repo{allsensor} bags carry wheel encoders} \textbar{} the baseline exists
  for static featureless-floor scenes only; no dynamic condition

\medskip
\noindent\textbf{40.}\quad \emph{table row:} E \textbar{} \flag{$n=5$ across four worlds, no condition repeated more than twice} \textbar{} no
  mean, deviation or interval; between-bag differences are not effects

\medskip
\noindent\textbf{41.}\quad \emph{table row:} E \textbar{} \flag{ORB-SLAM3 drops frames at a depth-5 best-effort subscriber; VINS buffers
  100 deep} \textbar{} input delivery is not matched between the two; no timing comparison
  is admissible

\medskip
\noindent\textbf{42.}\quad \emph{table row:} E \textbar{} \flag{ORB-SLAM3 run-to-run spread unmeasured, one run per configuration} \textbar{} loop-closure differences cannot be attributed; repeats required

\medskip
\noindent\textbf{43.}\quad \emph{table row:} E \textbar{} \flag{loop-closure state absent from \repo{run_metadata.csv}} \textbar{} configurations
  distinguishable only by inspecting \repo{loop_closing.csv}

\medskip
\noindent\textbf{44.}\quad \emph{table row:} E \textbar{} \flag{origin of the wheel yaw-rate offset untested} \textbar{} the drift is predicted
  but not yet attributable or correctable (\reportref{sec:h1-falsified})

\medskip
\noindent\textbf{45.}\quad \emph{table row:} E \textbar{} \flag{accelerometer characterisation unstable across bags} \textbar{} bias 0.019 to
  0.068, scale 1.020 to 1.067 for one simulated sensor

\medskip

\end{document}
